\documentclass[11pt]{article}
\usepackage[margin=1in]{geometry}
\usepackage{amsmath,amssymb}

\newcommand{\logx}{\log x}
\newcommand{\logwk}{\log(w,K,\tilde k)}

\title{Bnc Regime Neighbor Update Rule}
\author{}
\date{}

\begin{document}
\maketitle

\section{Inner Affine Map}

For a fixed Bnc regime, the fixed-point inner affine map is defined by
\[
    \logwk =
    A \logx + A_0,
    \qquad
    A =
    \begin{bmatrix}
        P_w\\
        N\\
        P\Pi
    \end{bmatrix},
    \qquad
    A_0 =
    \begin{bmatrix}
        P_{0,w}\\
        0\\
        P_0
    \end{bmatrix}.
\]
When \(A\) is nonsingular,
\[
    \logx = H_{\mathrm{inner}}\logwk + H_{0,\mathrm{inner}},
    \qquad
    H_{\mathrm{inner}}=A^{-1},
    \qquad
    H_{0,\mathrm{inner}}=-A^{-1}A_0.
\]

\section{One-Row Neighbor Update}

Every Bnc graph edge changes exactly one row of \(A\):
\[
    A' = A + e_s c,
    \qquad
    A_0' = A_0 + e_s c_0.
\]
Here \(s\) is the affected row of the inner system and \(c\log x+c_0=0\) is the crossing hyperplane in \(x\)-space.

Let
\[
    h = cH_{\mathrm{inner}},
    \qquad
    h_0 = cH_{0,\mathrm{inner}} + c_0,
    \qquad
    \alpha = \frac{1}{1+h_s}.
\]
If \(1+h_s\neq 0\), Sherman-Morrison gives
\[
    H_{\mathrm{inner}}'
    =
    H_{\mathrm{inner}}
    -
    \alpha\,
    H_{\mathrm{inner}}e_s\,h,
\]
and
\[
    H_{0,\mathrm{inner}}'
    =
    H_{0,\mathrm{inner}}
    -
    \alpha\,
    H_{\mathrm{inner}}e_s\,h_0.
\]
If \(1+h_s=0\), the target regime is singular with nullity one, and the limiting ray is
\[
    H_{\mathrm{inner}}' =
    -H_{\mathrm{inner}}e_s\,h,
    \qquad
    H_{0,\mathrm{inner}}' =
    -H_{\mathrm{inner}}e_s\,h_0.
\]

\section{Binding Edge}

For a binding graph edge, only the binding permutation changes. If the changed binding row is a \(q_{\mathrm{cat}}\) row, then neither \(P_w\) nor \(P\Pi\) changes, so
\[
    A'=A,\qquad H_{\mathrm{inner}}'=H_{\mathrm{inner}}.
\]
If the changed binding row is a \(w\) row, say row \(a\), then
\[
    s=a,
    \qquad
    c=e_{j_2}^{\top}-e_{j_1}^{\top},
    \qquad
    c_0=\log\frac{L_{a,j_2}}{L_{a,j_1}},
\]
after accounting for the stored edge orientation. The generic one-row update applies directly.

\section{Catalysis Edge}

A catalysis edge changes one row in the flux dominance matrix. Let its oriented flux-space hyperplane be
\[
    c_v\log v+c_{0,v}=0,
    \qquad
    \log v=\Pi\log x+\log k.
\]
If the edge changes the positive flux dominance of fixed-point equation \(a\), then
\[
    s=d_w+r+a,
    \qquad
    c=c_v\Pi,
    \qquad
    c_0=c_{0,v}.
\]
If it changes the negative flux dominance of fixed-point equation \(a\), then \(P=P^+-P^-\) changes with the opposite sign:
\[
    s=d_w+r+a,
    \qquad
    c=-c_v\Pi,
    \qquad
    c_0=-c_{0,v}.
\]
Again the same one-row update applies.

\section{Singular Direction Check}

For nullity-one targets, \(H_{\mathrm{inner}}\) is a ray. If two regular neighbors propagate into the same singular Bnc regime, the propagated matrices should be positive scalar multiples only when the determinant orientation is globally coherent. In mixed Bnc regimes this orientation is not guaranteed by the binding-network positivity argument, because the bottom block \(P\Pi\) is produced by catalysis dominance rather than conservation. The implementation therefore records when two propagated singular rays are not positive scalar multiples.

\end{document}
